\documentclass[UTF8,a4paper,11pt]{ctexart}
\usepackage[margin=22mm,headheight=14pt,footskip=12mm]{geometry}
\usepackage{amsmath,amssymb,array,longtable,booktabs,tikz,needspace,fancyhdr,hyperref}
\usepackage{microtype}
\hypersetup{colorlinks=false,pdfborder={0 0 0},pdfauthor={Good Learning},bookmarks=false}
\newfontfamily\circledfont{Segoe UI Symbol}
\newcolumntype{P}[1]{>{\raggedright\arraybackslash}p{#1}}
\setlength{\LTleft}{0pt}\setlength{\LTright}{0pt}
\setlength{\parindent}{0pt}\setlength{\parskip}{.65em}
\linespread{1.18}
\raggedbottom
\setcounter{secnumdepth}{0}
\ctexset{section={format=\Large\bfseries\sffamily,beforeskip=1.35em,afterskip=.65em},
subsection={format=\large\bfseries\sffamily,beforeskip=1em,afterskip=.45em}}
\newcommand{\chaptertitle}[1]{\Needspace{8\baselineskip}\noindent{\LARGE\bfseries\sffamily #1}\par\vspace{4mm}\hrule\vspace{6mm}}
\newcommand{\CompactQuestion}[2]{\par\noindent\textbf{#1.}\ #2\par\vspace{.8em}}
\newcommand{\PracticeQuestion}[3]{\par\Needspace{4\baselineskip}\noindent\begin{minipage}{\linewidth}\textbf{#1.}\ #2\par\noindent\rule{0pt}{#3mm}\par\end{minipage}\par}
\newcommand{\CompactQuestionColon}[2]{\par\noindent\textbf{#1：}\ #2\par\vspace{.8em}}
\newcommand{\PracticeQuestionColon}[3]{\par\Needspace{4\baselineskip}\noindent\begin{minipage}{\linewidth}\textbf{#1：}\ #2\par\noindent\rule{0pt}{#3mm}\par\end{minipage}\par}
\pagestyle{fancy}\fancyhf{}\fancyfoot[C]{\thepage}\renewcommand{\headrulewidth}{0pt}
\begin{document}
\thispagestyle{empty}
\noindent{\sffamily\footnotesize 备考路径}\hfill{\sffamily\footnotesize 学习导图}\par
\vspace{7mm}
\noindent{\sffamily\fontsize{25}{30}\selectfont\bfseries 人工智能备考计划}\par
\vspace{4mm}\noindent\rule{\linewidth}{.85pt}\par\vspace{6mm}
\begin{center}
\begin{tikzpicture}[x=1cm,y=1cm,>=stealth,every node/.style={font=\sffamily}]
\path[use as bounding box] (-8.1,-19.3) rectangle (8.1,.4);
\node[anchor=west,font=\sffamily\Large\bfseries] at (-7.5,-.55) {从章节练习到限时模拟};
\node[anchor=east,font=\sffamily\small,text=black!55] at (7.5,-.55) {11/15–12/08};
\draw[black!45,line width=.7pt] (-6.65,-1.75) -- (-6.65,-16.65);
\node[circle,draw=black,line width=.7pt,minimum size=9mm,fill=white,font=\sffamily\small\bfseries] at (-6.65,-2.25) {01};
\node[anchor=west,font=\sffamily\large\bfseries] at (-5.75,-1.8) {诊断与传统模型};
\node[anchor=east,font=\sffamily\small,text=black!55] at (7.45,-1.8) {11/15–11/22};
\node[anchor=north west,text width=12.8cm,align=left,font=\sffamily\small] at (-5.75,-2.4) {搜索、梯度、树分裂\quad{}·\quad{}验证三类主要步骤};
\draw[black!25,line width=.4pt] (-5.75,-3.6) -- (7.45,-3.6);
\node[circle,draw=black,line width=.7pt,minimum size=9mm,fill=white,font=\sffamily\small\bfseries] at (-6.65,-5.6) {02};
\node[anchor=west,font=\sffamily\large\bfseries] at (-5.75,-5.1499999999999995) {神经网络、视觉与语言};
\node[anchor=east,font=\sffamily\small,text=black!55] at (7.45,-5.1499999999999995) {11/23–11/29};
\node[anchor=north west,text width=12.8cm,align=left,font=\sffamily\small] at (-5.75,-5.75) {传播、卷积、注意力\quad{}·\quad{}逐步手算与复做};
\draw[black!25,line width=.4pt] (-5.75,-6.949999999999999) -- (7.45,-6.949999999999999);
\node[circle,draw=black,line width=.7pt,minimum size=9mm,fill=white,font=\sffamily\small\bfseries] at (-6.65,-8.95) {03};
\node[anchor=west,font=\sffamily\large\bfseries] at (-5.75,-8.5) {强化学习与跨章短测};
\node[anchor=east,font=\sffamily\small,text=black!55] at (7.45,-8.5) {11/30–12/02};
\node[anchor=north west,text width=12.8cm,align=left,font=\sffamily\small] at (-5.75,-9.1) {转移、贝尔曼、Q-learning\quad{}·\quad{}记录三个卡点};
\draw[black!25,line width=.4pt] (-5.75,-10.299999999999999) -- (7.45,-10.299999999999999);
\node[circle,draw=black,line width=.7pt,minimum size=9mm,fill=white,font=\sffamily\small\bfseries] at (-6.65,-12.3) {04};
\node[anchor=west,font=\sffamily\large\bfseries] at (-5.75,-11.850000000000001) {两次模拟与变式};
\node[anchor=east,font=\sffamily\small,text=black!55] at (7.45,-11.850000000000001) {12/03–12/06};
\node[anchor=north west,text width=12.8cm,align=left,font=\sffamily\small] at (-5.75,-12.450000000000001) {闭卷限时\quad{}·\quad{}用第二次检验第一次的错步};
\draw[black!25,line width=.4pt] (-5.75,-13.65) -- (7.45,-13.65);
\node[circle,draw=black,line width=.7pt,minimum size=9mm,fill=white,font=\sffamily\small\bfseries] at (-6.65,-15.65) {05};
\node[anchor=west,font=\sffamily\large\bfseries] at (-5.75,-15.200000000000001) {考前回忆};
\node[anchor=east,font=\sffamily\small,text=black!55] at (7.45,-15.200000000000001) {12/07–12/08};
\node[anchor=north west,text width=12.8cm,align=left,font=\sffamily\small] at (-5.75,-15.8) {三道错题的关键步骤\quad{}·\quad{}一页公式与条件};
\draw[black!25,line width=.4pt] (-5.75,-17.0) -- (7.45,-17.0);
\draw[black!45,line width=.7pt] (-7.5,-18.35) -- (7.5,-18.35);
\node[anchor=west,font=\sffamily\bfseries] at (-7.4,-18.8) {复盘};
\node[anchor=east,font=\sffamily\small] at (7.4,-18.8) {独立作答  →  定位错步  →  变式复做};
\end{tikzpicture}
\end{center}
\clearpage

\chaptertitle{人工智能备考计划}
\textbf{初版／时间与正式范围待核 · 2026 年 11 月 15 日—12 月 8 日}

12 月 8 日考试，考前可安排 23 天。11/15 和 12/07 各 45 分钟，11/22 和 11/29 各 60 分钟，12/06 为 105 分钟，其余日期每天 90 分钟，合计 \textbf{32 小时 15 分钟}，考试日可选回忆时间另计。学习时段可按当天课表移动。

“计算 7”“填空 42”等均指《人工智能复习资料.docx》各题型内部的原编号。K01–K21 对应《知识清单》。题库含五类题，下面按这五类练习，考试正式题型、时长和分值尚待核对。图、公式和表格直接看 Word 原件。

\section{阶段与验收}
\begin{itemize}
\item \textbf{11/15–11/22｜诊断、搜索与传统模型。} 起步抽测后，把路径、梯度和分裂题各做一遍。到 11/22 能独立写对其中至少两类的主要步骤，第三类若失手则加入 11/26 的回顾。
\item \textbf{11/23–11/29｜神经网络、视觉与 NLP。} 到 11/29 独立完成一道前向或反向传播、一道卷积或池化、一道注意力题，至少两道过程完整。剩余薄弱题在 12/04 变式复做。
\item \textbf{11/30–12/02｜强化学习与首轮收口。} 能从转移图列矩阵，独立写出两步分布和一种贝尔曼方程，再完成一次 Q 更新。12/02 做跨章短测，对无法独立完成的题标记到两次模拟之间订正。
\item \textbf{12/03–12/06｜两次题库模拟与订正。} 用未做题或来自用户材料的改写题，统计按时完题率与独立正确率。第二次重点检验第一次出现的错误；两次都错的步骤列入考前短清单。
\item \textbf{12/07–12/08｜考前回忆。} 12/07 完成三道旧错题的关键步骤复述，整理一页公式与适用条件，考前按实际开考时间决定是否做短回忆。
\end{itemize}
\section{11/15–11/22 · 搜索与传统模型}
\subsection{11/15 · 起点诊断｜45 分钟}
\textbf{K02、K08、K14、K17、K19。} 10 分钟打开题库并标计算题组。25 分钟闭卷抽做填空 7、22、50，计算 7 的第一步，以及计算 8 的一道相似度。最后 10 分钟对照资料，把五项记为“独立完成／有提示完成／待复做”，写下首个出错步骤和耗时。没学过的项直接标记即可。把后续每章首次开始前 5 分钟的回顾分给这些薄弱项。

\subsection{11/16 · 绪论与数学｜90 分钟}
\textbf{K01–K02。} 10 分钟回忆诊断，25 分钟过绪论、导数与极值、条件概率、矩阵规则，40 分钟做偏导、期望、矩阵乘法、行列式各一小题，15 分钟订正。验收为四题至少三题过程独立写全，并写出偏导涉及的变量和乘法维度。

\subsection{11/17 · 图搜索与最短路径｜90 分钟}
\textbf{K03–K04。} 10 分钟回顾队列与矩阵，25 分钟读搜索课件例题，40 分钟手推同一张小图的 DFS、BFS、A*，任选计算 16 或铁路网络的小子图做最短路，15 分钟对照。验收为访问次序、前沿记录与总成本，并注明边权和启发函数条件。

\subsection{11/18 · 局部搜索与博弈｜90 分钟}
\textbf{K03–K05。} 10 分钟重画昨日一轮前沿；25 分钟回顾双向搜索、爬山、退火、遗传与蒙特卡罗树搜索；40 分钟用黑白棋作业思路画一棵浅层博弈树，分别做极小极大和 Alpha–Beta；15 分钟核对剪枝。验收为根值、剪掉的子树及两条算法比较。

\subsection{11/19 · 数据划分、泛化与聚类｜90 分钟}
\textbf{K06–K07。} 10 分钟复述搜索条件，25 分钟复习划分、损失和泛化，40 分钟做简答 6–8 中两题及一轮 K 均值，口述谱聚类的相似图思路，15 分钟订正。验收为一张数据用途对照表与带新中心的分配表。

\subsection{11/20 · 梯度下降与线性模型｜90 分钟}
\textbf{K02、K06、K08。} 10 分钟重做错的一项泛化判断，25 分钟回顾目标函数、学习率、正则化与最大间隔，40 分钟做计算 7 及简答 4、9–11 中两题，15 分钟订正。验收为三步更新表，以及 L1/L2、回归/分类两组条件对照。

\subsection{11/21 · 决策树与集成｜90 分钟}
\textbf{K09–K10。} 10 分钟复述过拟合和平方误差，25 分钟练熵、加权分裂、叶预测、随机森林与 GBDT，40 分钟做计算 17 和填空 42–43，15 分钟核对。验收为两组候选分裂及叶值，写明准则和随机森林的两种随机性。

\subsection{11/22 · 第一轮复盘｜60 分钟}
\textbf{K03–K10。} 25 分钟闭卷做一道路径小题、一道梯度变式、计算 18 的分裂设置；20 分钟订正并用一组简单残差复述 GBDT 一轮；15 分钟补漏或休息。验收为三类核心题至少两类能独立写出主步骤，记录剩余卡点。

\section{11/23–11/29 · 神经网络、视觉与 NLP}
\subsection{11/23 · 激活函数与前向传播｜90 分钟}
\textbf{K08、K11。} 10 分钟复做一次梯度更新，25 分钟整理感知机、激活和多层网络，40 分钟做填空 32–35 及计算 2、19、20 中两题，15 分钟核对。验收为每层输入输出形状及 ReLU 处理结果。

\subsection{11/24 · 反传与 LSTM｜90 分钟}
\textbf{K11–K12。} 10 分钟闭卷写两层前向，25 分钟回顾链式法则、梯度消失和门控，40 分钟做计算 1、3 中一至两题及填空 29，15 分钟订正。验收为一次权重梯度计算与一张标明三道门的 LSTM 图。

\subsection{11/25 · 卷积、滤波与池化｜90 分钟}
\textbf{K12–K14。} 10 分钟回顾权值共享，25 分钟复习卷积输出尺寸、步长与填充，并过成像、边缘和视差，40 分钟做计算 4–6 中两题和填空 21–22，15 分钟订正。验收为两张特征图、池化结果与一句成像或视差条件解释。

\subsection{11/26 · 上采样、检测与分割｜90 分钟}
\textbf{K09、K13–K15。} 10 分钟重算一个卷积尺寸或补 11/22 的薄弱题，25 分钟回顾双线性插值、转置卷积和 Faster R-CNN，40 分钟做计算 9、22–24 中三小题，15 分钟对照。验收为一题上采样手算、检测与分割对比以及 RPN 分类和回归的输出数量。

\subsection{11/27 · 语言模型、词向量与解码｜90 分钟}
\textbf{K16–K17。} 10 分钟回顾条件概率和内积，25 分钟过分词、n-gram、困惑度、CBOW/Skip-Gram，40 分钟做填空 14–15、27–28 和计算 8、25 中的代表小问，15 分钟订正。验收为一张概率或解码表、两组余弦值和两种词嵌入预测方向。

\subsection{11/28 · 注意力与位置编码｜90 分钟}
\textbf{K11、K17–K18。} 10 分钟重算一道相似度，25 分钟梳理 Seq2Seq、Q/K/V、多头、位置和因果掩码，40 分钟做计算 10 或 27 及填空 47–48，15 分钟核对。验收为写出一个查询完整的投影、缩放、归一化和加权步骤。

\subsection{11/29 · 第二轮复盘｜60 分钟}
\textbf{K11、K14、K18。} 30 分钟从计算 3、6、27 的剩余原题或改写题抽三个短问，20 分钟逐步订正，10 分钟补漏或休息。验收为三题至少两题过程独立完整，给失手项约定 12/04 的变式。

\section{11/30–12/02 · 强化学习与跨章短测}
\subsection{11/30 · 马尔可夫链｜90 分钟}
\textbf{K02、K19。} 10 分钟回顾矩阵与期望，25 分钟区分无后效性、齐次性、稳态与吸收状态，40 分钟做计算 12–15 中两题，15 分钟订正。验收为一份两步转移结果、一个稳态方程和归一性检查。

\Needspace{7\baselineskip}
\subsection{12/01 · 贝尔曼方程与 Q-learning｜90 分钟}
\textbf{K19–K20。} 10 分钟回顾两步转移，25 分钟整理 MDP、折扣回报、两种价值函数、值迭代，40 分钟做计算 11、填空 50 或 51 及一道 Q-learning 单步更新题（新编或课件例题），15 分钟订正。验收为两状态的贝尔曼方程、一条 Q 更新与终止状态处理。

\subsection{12/02 · 第一轮收口｜90 分钟}
\textbf{K01–K21。} 15 分钟闭卷列章节与公式条件，45 分钟用未做题快测，每类题至少一项，覆盖搜索、树、神经网络、视觉、NLP、强化学习，20 分钟标出前三个卡点，10 分钟核对补充范围 K21。验收为有耗时和首个错步的记录，以及模拟前应重做的三个步骤。

\section{12/03–12/06 · 模拟与变式}
\subsection{12/03 · 题库模拟一｜90 分钟}
\textbf{K03–K20。} 60 分钟闭卷做 6 判断、6 填空、3 名词解释、2 简答、3 个计算小问。计算优先搜索或分裂、视觉或深度网络、强化学习。30 分钟按步骤订正。验收为按时完题率、独立正确率、计算题错步和下一次要替换的题目。题目优先选未做过的原题。

\subsection{12/04 · 攻克两到三个错步｜90 分钟}
\textbf{K03–K20，依模拟一结果选取。} 15 分钟闭卷重现错步，25 分钟复查对应课件，35 分钟做两到三个变式，15 分钟自检。验收为每个错步都能写一条具体检查规则，并记录变式是否独立做对。

\subsection{12/05 · 混合回忆与保留题｜90 分钟}
\textbf{K06、K09–K12、K14–K20。} 15 分钟回忆公式和条件，45 分钟做四个跨章小问（默认梯度或分裂、Q/K/V、转置卷积或 RPN、贝尔曼或 Q 更新），20 分钟订正，10 分钟挑选模拟二的未见题。验收为四题的过程与易混条件，保留模拟二题目至次日开做。

\subsection{12/06 · 题库模拟二｜105 分钟}
\textbf{K03–K20。} 70 分钟闭卷做 6 判断、6 填空、3 名词解释、3 简答、4 个计算小问。用未见题或标为“来自用户材料的改写”的变式，其中至少一题验证模拟一薄弱项。35 分钟订正。验收为两次完题率和独立正确率的对照、复错题及最后三项复习顺序。

\section{12/07–12/08 · 考前回忆}
\subsection{12/07 · 最后一页｜45 分钟}
\textbf{K04、K08–K11、K14、K18–K20 中尚不熟的项。} 15 分钟默写公式、条件与输入输出形状，20 分钟闭卷复做三道旧错题的关键步骤，10 分钟核对考试时间、地点及携带要求。验收为一页回忆清单，重点留容易漏写的条件和解题步骤。

\subsection{12/08 · 考试日｜可选 10–15 分钟}
按实际开考时间回忆前一天清单上的三个典型错步，随后按考试安排到场。验收为能口述每条检查规则。

\section{如何按结果调整}
\begin{itemize}
\item 每次记录来源和题号、耗时、是否独立做对、首个错步。完题率用限时内做完的题数除以总题数，独立正确率用独立答对的题数除以已尝试题数，过程题缺关键步骤记为待复做。
\item 普通复习日优先用开头 10 分钟复做近一两天的错题，阶段复盘再抽一题。两次独立完成后，把该类移入下一次混合回忆。重复卡在同一步时，次日挪 15–20 分钟做一个例题和一个变式。
\item 若某天只有 45 分钟，用 5 分钟回忆、25 分钟攻一类典型题、15 分钟订正。积压的题优先移到 11/22、11/29、12/04 的补漏时段；先减少 K01、K21 与重复题量，保留两次模拟、订正和考前回忆。
\item 拿到教师划定范围、考试时长或样卷后，更新 K 编号对应范围和两套模拟的时间、题型配比。题库中缺图或参数冲突的题先核对 Word 原件与课件；仍有歧义时按不同给定条件练习，记录待确认之处。
\item 依据为《人工智能复习资料.docx》、课件和课程作业。现有复习资料未提供正式分值、考试时长或标准答案。笔记覆盖多章，实际掌握程度以作答验证。初步推荐“复习”模式；若多个章节的基础例题尚无法起步，可从这些章节切换为“零基础”讲解。
\end{itemize}
\end{document}
